\documentclass[10pt,a4paper]{amsart}
\usepackage[margin=19mm]{geometry}
\usepackage{amsmath,amssymb,mathtools}
\usepackage[scr=rsfs,cal=euler]{mathalfa}
\usepackage{xfrac}
\usepackage[unicode,hidelinks,bookmarks=false]{hyperref}
\newcommand{\ie}{i.e.\ }
\newcommand{\cf}{cf.\ }
\newcommand{\resp}{respectively\ }
\newcommand{\Subsection}{\subsection}
\providecommand{\nobreakdash}{\nobreak\hbox{-}}
\setlength{\emergencystretch}{2em}
\multlinegap0pt
\usepackage{cancel}
\usepackage{enumerate}
\usepackage{stackengine}
\usepackage{amssymb}
\usepackage{bm}
\usepackage{dsfont}
\usepackage{mdframed}
\usepackage{xcolor}
\usepackage{mathtools}
\usepackage{bbm}
\usepackage{tikz}
\usepackage[normalem]{ulem}
\newtheorem{theorem}{Theorem}
\newtheorem{proposition}{Proposition}
\newtheorem{lemma}{Lemma}
\usepackage{cleveref}
\crefname{theorem}{Theorem}{Theorems}
\crefname{proposition}{Proposition}{Propositions}
\crefname{lemma}{Lemma}{Lemmas}
\newcommand{\N}{\mathbb{N}}
\newcommand{\R}{\mathbb{R}}
\newcommand{\T}{\mathbb{T}}
\newcommand{\Z}{\mathbb{Z}}
\newcommand{\bS}{\mathbb{S}}
\newcommand{\ind}[1]{\mathbbm{1}_{#1}}
\title{{\bfseries Structure and Spectrum of Nonergodic Nilsystems}}
\author{By~~{\scshape Felipe~Hernández~}}
\date{Source: arXiv:2602.21021v2 (CC BY 4.0)}
\begin{document}
\maketitle

\noindent\emph{Source and licence.} {\bfseries Structure and Spectrum of Nonergodic Nilsystems}. Version arXiv:2602.21021v2 (CC BY 4.0), distributed by arXiv under the Creative Commons Attribution 4.0 International licence, http://creativecommons.org/licenses/by/4.0/, as recorded in the arXiv licence metadata for this exact version and shown on the abstract page https://arxiv.org/abs/2602.21021v2. Full licence notice: \texttt{licenses/arxiv-2602.21021-CC-BY-4.0.txt}.\par\smallskip

\noindent\textit{Editorial scope.} Counted unit: the article's theorem Theo-B (source0.tex lines 267--271). The article's complete original proof of it is delivered with it (source0.tex lines 644--675, one \texttt{proof} environment of 32 lines, which the article places away from the statement). No mathematics is written, repaired or re-derived: the statement and the proof are the article's own lines, in the article's own notation, and no retained line is substituted or re-flowed.  Retained uncounted as the source-local context the proof needs: lines 330--337; lines 341--343; lines 492--494; lines 520--522; lines 573--575; lines 591--594; lines 628--631.  Nothing was omitted from any retained span, so the package carries no omission record.  No retained line contains a citation, so the wrapper carries no bibliography and no bibliography entry.  No retained line contains a figure environment, an image-inclusion command or drawing code, so no image is delivered and the package declares no asset dependency.  Editorial formatting only, in addition: an AMS \texttt{amsart} preamble that restates the article's own declarations and macros used by the retained lines, namely the environments \texttt{theorem}, \texttt{proposition}, \texttt{lemma} on separate counters, together with the standard packages those lines need.  The retained environments print in this document as Theorem 1, Proposition 1, Lemma 1 in order of appearance, whereas the article prints the same items with the numbering of its own full document.  The complete native source of the article as distributed in the arXiv e-print for this exact version is bundled unchanged as \texttt{sources/195236/source0.tex}.
\begin{proposition}\label{PW}
    Let $X=G/\Gamma$ be a nilmanifold. Let $L\leq Z(G)$ be a rational group and $J=L/(L\cap \Gamma)$ be the associated compact abelian group. Then, we have the decomposition 
    $$L^2(X)=\bigoplus_{\chi\in \widehat{J}}V_\chi, $$
    where $\widehat{J}$ is the character group of $J$ and for $\chi\in \widehat{J}$ we denote 
    \begin{equation}\label{def-V-chi}
       V_\chi:=\{ f\in L^2(X) : f(gx)= \chi(g)f(x), \forall g\in L, \text{ for }\mu\text{-a.e. } x\in X\}.  
    \end{equation}
\end{proposition}

\begin{proposition}\label{Lemma-measure-in-circle}
    Let $\sigma$ be a Borel measure on $\bS^1$ such that for each $\alpha\in \bS^1$, $\sigma^\alpha\ll \sigma$. Then $\sigma\thicksim m_{\bS^1}$.
\end{proposition}

\begin{theorem}\label{Spectrum-uniform-part}
    Let $(X=G/\Gamma,\mu,T)$ be a nilsystem and $H\subseteq G$ its Leibman group. Then $T$ exhibits infinite Lebesgue spectrum in $L^2([H,H] \backslash X)^\perp$.
\end{theorem}

\begin{theorem}\label{main-theorem}
    Let $(X=G/\Gamma,\mu,T)$ a nilsystem. Then $T$ exhibits pure discrete spectrum on $L^2(\mathcal{J}([\tau,G],\Gamma)\backslash X)$. 
\end{theorem}

\begin{proposition}\label{app-coarea-formula}
    Let $f\in \mathcal{C}^1(\R^d,\R)$ be such that the set $\{x\in \R^d : \nabla f (x)=0\}$ has zero Lebesgue measure. Then the pushforward $f^*(m_{\R^d})$ of the Lebesgue measure is absolutely continuous with respect the Lebesgue measure $m_{\R}$ on $\R$.
\end{proposition}

\begin{proposition}\label{tj-explicit}
    Let $(X=G/\Gamma,\mu,T=T_\tau)$ be a nilsystem of order $1$, and $H$ its Leibman group. Let $Y=G/H\Gamma$ be the factor defined by the invariant sigma algebra $\mathcal{I}(T)$ of $X$ and $(\mu_y)_{y\in Y}$ the ergodic desintegration given by it. Let $\eta_{G}: \R^d \to G^\circ$ be the Mal'cev coordinates of $G^\circ$ adapted to $H^\circ$. Let $(\chi_j)_{j\in \N}\subseteq \widehat{H}$ be an orthonormal basis of $L^2(H/H\cap \Gamma)$. Then, there is $l\leq d$ such that for each $j\in \N$ the function $t_j:Y\to \T$ defined for $\mu_Y$-a.e. $y=gH\Gamma$ with $g\in \eta_G(\{0\}^{d-l}\times [0,1)^l)$, by $$e(t_j(y))= \chi_j(g^{-1}\tau g)$$
    is a Borel map such that  for $\mu_Y$-a.e. $y\in Y$, $\mu_y(\{e(t_j(y))\}_{j\in \N})=1$. 
\end{proposition}

\begin{lemma}\label{relative-basis}
    Let $(X,\mu,T)$ be a system of order one and $f\in L^2(\mu)$. Let $(Y,\mathcal{Y},\nu,T)$ be the factor defined by the invariant sigma algebra $\mathcal{I}(T)$, where $(\mu_y)_{y\in \N}$ denotes the ergodic decomposition given by this factor. If there are Borel maps $t_j:Y\to \T$ for $j\in \N$, such that for $\mu_Y$-a.e. $y\in Y$, the Haar measure $\mu_y$ on $\overline{\mathrm{Orb}}_{T}(y)$ is supported in $\{t_j(y)\}_{j\in \N}$, then there exists a relative orthonormal basis of eigenfunctions $(\phi_j)_{j\in\N}$ with eigenvalues $\lambda_j(y) \in\{\ind{A_i}(y) e(t_i(y))\}_{i\in \N}$ where 
    $$A_i=\{y\in Y : \mu_y(\{e(t_i(y))\})>0 \}. $$
\end{lemma}

\begin{theorem}\label{Theo-B}
Let $(X=G/\Gamma,\mu,T_{\tau})$ be a nilsystem. Then the Koopman representation of $T_{\tau}$ on $L^2(X,\mu)$ decomposes as
  $$ L^2(X)=L^2(\mathcal{J}([\tau,G],\Gamma)\backslash X)\oplus L^2(\mathcal{J}([\tau,G],\Gamma)\backslash X)^\perp,$$
  where $T_{\tau}$ exhibits discrete spectrum on $L^2(\mathcal{J}([\tau,G],\Gamma)\backslash X)$ and infinite Lebesgue spectrum on $L^2(\mathcal{J}([\tau,G],\Gamma)\backslash X)^\perp$ if nontrivial. In particular, $G/\mathcal{J}([\tau,G],\Gamma)\Gamma$ is the discrete-spectrum factor of $X$.
\end{theorem}

\begin{proof}[Proof of \cref{Theo-B}]
Let $H$ be the Leibman group of $(X=G/\Gamma,\mu,T)$. First, we notice that if $[\tau,G]=\{e_G\}$ then $L^2(X)=L^2(\mathcal{J}([\tau,G],\Gamma)\backslash X)$ and the result follows from \cref{main-theorem}. Thus, we will assume that $[\tau,G]\neq \{e_G\}$. Second, by \cref{Spectrum-uniform-part}, $L^2([H,H]\backslash X)^\perp$ exhibits infinite Lebesgue spectrum. Hence, it is enough to prove the result in $L^2([H,H]\backslash X)$. Then, replacing $X$ by $[H,H]\backslash X$, we assume without loss of generality that $[H,H]=\{e_G\}$ and thus $H$ is an abelian.


Let $(\chi_j)_{j\in \N}$ be a set of characters in $\widehat{H}$ forming an orthonormal basis of $L^2(H/(H\cap \Gamma))$, such that $\chi_j(\Gamma\cap H)\equiv 1$. Denote by $F= \{0\}^{d-l} \times [0,1)^l$ the fundamental domain for the projection of $G^\circ/H^\circ$ onto $Y$. Using  \cref{relative-basis}  with the family of functions $\{t_j\}_{j\in \N}$ given by \cref{tj-explicit}, we obtain the decomposition
    $$L^2(X)= \bigoplus_{j\in \N} \overline{L^{\infty}(X,\mathcal{I}(T)) \phi_j  }, $$
    where for each $j\in \N$, $\phi_j$ has eigenvalue $\lambda_j(y)=  \ind{A_{i_j}}(y) e(t_{i_j}(y))$ where $A_{i_j}=\{y \in Y : \mu_y(\{e(t_{i_j}(y))\})>0 \}.$ For the sake of notation, let us denote $i_j=j$. Thus, it is enough to see that the spaces $E_j:=  \overline{ L^\infty(X,I(T))  \phi_j} $ have maximal spectral type absolutely continuous with respect to the Lebesgue measure on $\bS^1$, whenever $\phi_j$ is not $\mathcal{J}([\tau,G],\Gamma)$-invariant. Let $\psi \in L^\infty(X,I(T))$ and $j \in \N$ be such that $\psi\phi_j\not\equiv 0$ and $\phi_j$ are not $\mathcal{J}([\tau,G],\Gamma)$-invariant. We have 
$$\int_X (\overline{\psi\phi_j})\cdot T^n(\psi\phi_j) d\mu=\int_Y |\psi \phi_j|^2 \ind{A_j} e(t_j(y))^nd\mu_Y. $$
In this way, the spectral type of $f= \psi\phi_j$ is given by the pushforward of $|\psi \phi_j|^2\ind{A_j}\mu_Y  $  through the map $y\in Y \mapsto e(t_j(y))\in \bS^1$. We claim $|\psi \phi_j|^2 \ind{A_j}\not\equiv 0$. Indeed, otherwise $T(\psi\phi_j)= \psi \phi_j \ind{A_j} \lambda_j=0$ and thus $\psi\phi_j\equiv 0$ which is a contradiction. It follows that it is enough to prove that the measure $\nu$ defined by the Fourier coefficients
$$\hat{\nu}(n):=\int_Y e(t_j(y))^n d\mu_Y  $$
is absolutely continuous with respect to Lebesgue. Let $\eta_{H}:\R^{d-l}\to H^\circ, \text{ and }\eta_{G}: \R^d \to G^\circ$, be Mal'cev coordinates on $H^\circ$ and $G^\circ$ respectively. Since $\chi_j$ is a character in $H$, its restriction to $H^\circ$ must have the form
$$\chi_j|_{H^\circ}(h)= e_j( k_j^T \cdot \eta_H^{-1}(h) ) ,$$
for each $h\in H^\circ$, and for some $k_j\in \Z^{d-l}$. By definition of $t_j$, for $x\in [0,1)^l$ and $g=\eta_G(x)$ we have
$$e(t_j(y))= \chi_j(\tau) \chi_j([g^{-1},\tau] )= \chi_j(\tau) e( k_j^T \cdot \eta_H^{-1}([g^{-1},\tau] ))=\chi_j(\tau) e( k_j^T \cdot \eta_H^{-1}([\eta_G(x)^{-1},\tau] )). $$
Since the operation in $G$ is polynomial in $\R^d$, the function
$$x\in \R^d \to k_j^T \cdot \eta_H^{-1}([\eta_G(x)^{-1},\tau] ) $$
is a polynomial, which we denote by $p_j(x)$. Notice that $\chi_j$ is not $[\tau,G]$-invariant. Indeed, if for the sake of contradiction we assume that $\chi_j$ is $[\tau,G]$-invariant, then we see that $y\mapsto e(t_j(y))$ is constant. This implies that the function $y\mapsto \lambda_j(y)$ has value $0$ or $1$. By our definition of eigenfunction, we have that $\phi_j$ is $T$-invariant, and thus $H$-invariant. Since $[\tau,G]\subseteq H$, we have that $\mathcal{J}([\tau,G],\Gamma)\subseteq H$  as $H$ is a rational connected normal subgroup of $G^\circ$ containing $[\tau,G]$, and $\mathcal{J}([\tau,G],\Gamma)$ is the smallest with such a property. As $\phi_j$ is not $\mathcal{J}([\tau,G],\Gamma)$-invariant, this is a contradiction. Thus, $\chi_j$ is not $[\tau,G]$-invariant, which implies that $p_j$ is not constant. 
We have  
\begin{align*}
  \widehat{\nu}(n)&= \chi(\tau)^n\int_{[0,1)^l} \chi([\eta_G(x)^{-1},\tau] H\cap \Gamma)^n dm_{\R^d}(x)=\chi(\tau)^n \int_{[0,1)^l} e(p_j(x))^n dm_{\R^d}(x).   
\end{align*}
Thus, it is enough to show that the pushforward $p(m_{\R^d})$ is absolutely continuous with respect to the Lebesgue measure of $\R^d$. Since $p$ is not constant, then $\nabla p(x)\neq 0$ for $\mu$-a.e. $x\in \R^d$. By \cref{app-coarea-formula} we have that the push forward of $p_j^*(m_{\R^d})$ is absolutely continuous with respect to the Lebesgue measure, concluding that $ \nu \ll m_{\bS^1}$ and thus $\sigma_f\ll m_{\bS^1} $.

To see that the maximal spectral type is countable Lebesgue, it is enough to work in a factor of the system. Let $s\geq 2$ be the minimum such that the projection of $\mathcal{J}([\tau,G],\Gamma)$ in $G/G_{s}$ is non trivial. Reducing to the factor $G/G_s\Gamma$ of $X$, we can assume without loss of generality that $X=G/\Gamma$ is a $(s-1)$-step nilsystem such that $[\tau,g]\in Z(G)$ for each $g\in G$. Using \cref{PW} with $L=Z(G)$, we obtain the decomposition 
$$L^2(X)=\bigoplus_{\chi\in \widehat{J}}V_\chi, $$
    where $J= L/L\cap \Gamma$ and for $\chi\in \widehat{J}$ with $\chi(L\cap \Gamma)\equiv 1$, $V_\chi$ is defined as in \cref{def-V-chi}. 
    Since $\mathcal{J}([\tau,G],\Gamma)\subseteq Z(G)$ is nontrivial and connected, there are infinitely many $\chi\in \hat{J}$ such that $\chi$ is nontrivial in $\mathcal{J}([\tau,G],\Gamma)$. Take one of such $\chi$ and observe that $\chi([\tau,G])=\bS^1$ by connectedness of $[\tau,G]$ and continuity of $\chi$. Let $g\in G$ and let $f\in V_\chi$ attaining the maximal spectral type $\sigma$ on $V_\chi$. Let $g\in G$ and denote $\alpha:=  \chi([\tau,g])$. Similarly to the proof of \cref{Spectrum-uniform-part}, we have that 
    \begin{equation*}
       \int  \overline{T_gf} \cdot T_\tau^nT_gfd\mu= \int_{\bS^1} z^n d\sigma^\alpha(z),
    \end{equation*}
where we recall that $\sigma^\alpha$ denotes the measure $\sigma$ translated by $\alpha$. Thus, $\sigma^\alpha\ll \sigma$ for all $\alpha\in \chi([\tau,G])=\bS^1    $. By \cref{Lemma-measure-in-circle} we conclude that $\sigma=m_{\bS^1}$, concluding that in each subspace $V_\chi$ the maximal spectral type is the Lebesgue measure $m_{\bS^1}$, finishing the proof.  
\end{proof}

\end{document}
